\begin{table*}[t]
\caption{空间提示、反馈和融合的消融实验。前三项评价三维推理之前的二维引导，最后一项评价最终三维输出的固定视角投影。前两行为比较基线，随后四行为相对完整方法的组件消融，最后一行为完整方法；全表尚无测量结果。}
\label{tab:ablation}
\centering\small
\setlength{\tabcolsep}{7pt}
\begin{tabularx}{\textwidth}{@{}X c c c c@{}}
\toprule
配置 & \shortstack{二维匹配 mIoU\,(\%) $\uparrow$} & \shortstack{二维 BF1\,(\%) $\uparrow$} & \shortstack{二维部件 F1@0.5\,(\%) $\uparrow$} & \shortstack{投影 mIoU\,(\%) $\uparrow$}\\
\midrule
Automatic SAM2 & TBD & TBD & TBD & TBD\\
Qwen--SAM2（框＋点） & TBD & TBD & TBD & TBD\\
本文方法：仅单个正点 & TBD & TBD & TBD & TBD\\
本文方法：移除负点 & TBD & TBD & TBD & TBD\\
本文方法：移除反馈修正 & TBD & TBD & TBD & TBD\\
本文方法：移除点感知融合 & TBD & TBD & TBD & TBD\\
本文完整方法 & TBD & TBD & TBD & TBD\\
\bottomrule
\end{tabularx}
\vspace{2pt}
\begin{minipage}{\textwidth}\footnotesize
注：“仅单个正点”保留框、负点、反馈及融合，其与独立框＋点基线的差异必须由配置清单说明。移除点感知融合时的替代规则尚需固定；若同时删除多个操作，结果只能解释为该组合模块的贡献。匹配基于区域重叠，不按调色板 RGB 值一一对应。
\end{minipage}
\end{table*}
